\chapter{Finite perimeter and reduced boundary}\label{sec:finite-perimeter}
A topological boundary can change drastically when a null set is added. The distributional derivative of the characteristic function provides a boundary notion that depends only on the measurable set. Fleming describes the development of this approach in \cite{Fle16}, Section 3, pp.~3--4. The precise statements below follow \cite{Sim18}, Chapter 2, Section 2, and Chapter 3, Section 4.

\begin{definition}[Locally bounded variation]\label{def:bv}
For $u\in L^1_{\mathrm{loc}}(U)$ and an open set $W\subset U$, put
\[
 |Du|(W)=\sup\left\{\int_W u\,\operatorname{div}X\,d\mathcal{L}^n:
 X\in C_c^1(W;\mathbb{R}^n),\ |X|\leq1\right\}.
\]
We say $u\in BV_{\mathrm{loc}}(U)$ if this quantity is finite for every $W\Subset U$. Equivalently, its distributional derivative $Du$ is a vector-valued Radon measure, characterized by
\[
 \int_U u\,\operatorname{div}X\,d\mathcal{L}^n=-\int_U X\cdot dDu.
\]
The notation $|Du|$ then denotes its total variation measure. Proposition \ref{prop:bv-vector-representation} proves this equivalence and identifies the measure with the displayed supremum. We write $u\in BV(U)$ when $u\in L^1(U)$ and $|Du|(U)<\infty$. See \cite{Sim18}, Chapter 2, Section 2, pp.~52--53.
\end{definition}

\begin{definition}[Finite perimeter]\label{def:finite-perimeter}
A Lebesgue measurable set $E\subset\mathbb{R}^n$ has \emph{locally finite perimeter in $U$} if $\chi_E\in BV_{\mathrm{loc}}(U)$. Its perimeter in an open set $W\subset U$ is
\[
 P(E;W)=|D\chi_E|(W).
\]
It has finite perimeter in $U$ if $P(E;U)<\infty$; this does not require $\mathcal{L}^n(E\cap U)<\infty$. Sets equal almost everywhere have the same distributional derivative and perimeter. See \cite{Sim18}, Chapter 3, p.~96.
\end{definition}

\begin{proposition}[Lower semicontinuity of perimeter]\label{prop:perimeter-lsc}
Suppose $\chi_{E_j}\to\chi_E$ in $L^1_{\mathrm{loc}}(U)$. For every open $W\subset U$,
\[
 P(E;W)\leq\liminf_{j\to\infty}P(E_j;W),
\]
where the variational definition permits value $+\infty$ even before local finite perimeter is known.
\end{proposition}
\begin{proof}
For every $X\in C_c^1(W;\mathbb{R}^n)$ with $|X|\leq1$, local $L^1$ convergence and boundedness of $\operatorname{div}X$ give
\[
 \int_W\chi_E\operatorname{div}X\,d\mathcal{L}^n
 =\lim_j\int_W\chi_{E_j}\operatorname{div}X\,d\mathcal{L}^n
 \leq\liminf_j P(E_j;W).
\]
Take the supremum over $X$. This is the variational lower semicontinuity used in \cite{Sim18}, Chapter 2, Section 2.
\end{proof}

\begin{definition}[Reduced boundary and outward normal]\label{def:reduced-boundary}
Let $E$ have locally finite perimeter in $U$. Its \emph{reduced boundary} $\partial^*E\cap U$ consists of points $x\in U\cap\operatorname{supp}|D\chi_E|$ for which
\[
 v_E(x)=\lim_{r\downarrow0}\frac{D\chi_E(B_r(x))}{|D\chi_E|(B_r(x))}
\]
exists and has length one. We define the outward measure-theoretic unit normal by $\nu_E(x)=-v_E(x)$. For small $r$ the balls lie in $U$; membership in the support ensures that their denominators are positive. This is \cite{Sim18}, Chapter 3, equations (4.2)--(4.3), pp.~96--97, with the sign reversed to use the outward normal.
\end{definition}

\begin{theorem}[De Giorgi structure and Gauss--Green]\label{thm:de-giorgi}
Let $n\geq2$ and let $E$ have locally finite perimeter in $U\subset\mathbb{R}^n$. Then $\partial^*E\cap U$ is countably $(n-1)$-rectifiable and
\[
 |D\chi_E|=\mathcal{H}^{n-1}\llcorner(\partial^*E\cap U),
 \qquad D\chi_E=-\nu_E\,\mathcal{H}^{n-1}\llcorner(\partial^*E\cap U).
\]
Consequently, for every $X\in C_c^1(U;\mathbb{R}^n)$,
\[
 \int_E\operatorname{div}X\,d\mathcal{L}^n
 =\int_{\partial^*E\cap U}X\cdot\nu_E\,d\mathcal{H}^{n-1}.
\]
Here $X$ is extended by zero outside $U$. The boundary measures in these identities are taken inside $U$.

The proof below separates blow-up, graph covering, and measure identification. See \cite{Sim18}, Chapter 3, Theorem 4.4, pp.~97--100. The Gauss--Green identity follows immediately from the measure representation and Definition~\ref{def:bv}. The historical account in \cite{Fle16}, Section 3, uses the same outward-normal convention as this entry.
\end{theorem}
\begin{proof}
Theorem \ref{thm:reduced-boundary-graph-cover} proves countable rectifiability from the full pointwise blow-up and density estimates. Theorem \ref{thm:perimeter-hausdorff-identification} identifies the variation measure with Hausdorff measure on the reduced boundary. Its proof explicitly uses Theorem \ref{thm:radon-function-differentiation}, the abstract Radon--Nikodym theorem, and the cited rectifiable-set density theorem; the Besicovitch selection underlying Radon differentiation, the abstract representation theorem and the rectifiable-set density theorem remain explicit external inputs.

To identify the vector derivative, write $D\chi_E=-\nu\,|D\chi_E|$ using its measurable unit polar field. The differentiation argument in Corollary \ref{cor:perimeter-concentrated-reduced-boundary} gives $\nu=\nu_E$ at almost every point for the perimeter measure. Substituting the scalar measure representation proves the second displayed identity. Finally Definition \ref{def:bv} yields $\int_E\operatorname{div}X=-\int X\cdot dD\chi_E$, which is the stated Gauss--Green formula with outward normal. This conclusion depends on the measurable set modulo null sets; no regularity of its chosen topological boundary is required.
\end{proof}

\begin{example}[Smooth domains]\label{ex:smooth-perimeter}
Let $E\subset\mathbb{R}^n$ be a bounded open set with compact $C^1$ boundary. The classical divergence theorem identifies
\[
 D\chi_E=-\nu_E\,\mathcal{H}^{n-1}\llcorner\partial E,
 \qquad P(E;U)=\mathcal{H}^{n-1}(\partial E\cap U).
\]
The continuous outward normal has unit length, so the variation of the displayed vector measure is exactly the displayed surface measure. Thus the variational perimeter agrees with the usual boundary area. This also checks the sign in Theorem~\ref{thm:de-giorgi}.
\end{example}
